\chapter{有限作用时间如何选择电子的出射通道}
\label{chap:feqo-transition-kernel}

电子态与电磁场已经有了各自的动力学。让二者相互作用以后，实验首先问电子可能以哪些能量和动量离开。若直接使用无限长时间的能量守恒 \(\delta\) 函数，便无法描述短脉冲的谱宽；这里从有限时间演化开始计算。

\begin{definition}[相互作用绘景与精确传播]
将总哈密顿量分为可单独求解的 \(H_0\) 和相互作用 \(V(t)\)。在 \(H_0\) 的本征态中，\(V_I(t)=e^{\ii H_0t/\hbar}V(t)e^{-\ii H_0t/\hbar}\)。相互作用绘景的传播算符 \(U_I(t,t_0)\) 由
\begin{equation}
\ii\hbar\partial_t U_I(t,t_0)=V_I(t)U_I(t,t_0),\qquad U_I(t_0,t_0)=I
\label{eq:feqo-interaction-evolution}
\end{equation}
唯一确定。它的形式精确解是时间有序指数 \(\mathcal T\exp[-(\ii/\hbar)\int_{t_0}^{t}V_I(s)\dd s]\)；时间排序表示较晚时刻的算符写在左边。
\end{definition}

把式~\eqref{eq:feqo-interaction-evolution} 积分一次并在右边递推，得到 Dyson 级数。其一阶项给出有限时谱线；只有当后续项相对它足够小时，才把一阶概率当成实验预测。

\section{从时间积分到可分辨的能量宽度}
\label{chap:feqo-finite-time}

取 \(H_0|i\rangle=E_i|i\rangle\)、\(H_0|f\rangle=E_f|f\rangle\)，两态都已归一化。设吸收项在薛定谔绘景中为 \(V_{fi}w(t)e^{-\ii\omega t}\)，其中 \(V_{fi}\) 是能量量纲的矩阵元，实包络 \(w\) 无量纲。由相互作用绘景的定义，这个矩阵元还乘 \(e^{\ii(E_f-E_i)t/\hbar}\)。于是从 \(t=-\infty\) 到 \(+\infty\) 的一阶振幅为
\begin{equation}
c_f^{(1)}=-\frac{\ii V_{fi}}{\hbar}\int_{-\infty}^{\infty}\dd t\,w(t)e^{\ii\Delta_{fi}t},\qquad
\Delta_{fi}=\frac{E_f-E_i}{\hbar}-\omega .
\label{eq:feqo-finite-time-amplitude}
\end{equation}
失谐 \(\Delta_{fi}\) 的单位是 \(\si{\per\second}\)。它为零时，扰动的振荡与两态的能量差同步。

这一阶式来自一个有明确余项的演化方程。把式~\eqref{eq:feqo-interaction-evolution} 从 \(t_0\) 积到 \(t\)，再把所得的 \(U_I\) 代回积分号内，便有
\[
U_I(t,t_0)=I-\frac{\ii}{\hbar}\int_{t_0}^t V_I(t_1)\dd t_1
-\frac1{\hbar^2}\int_{t_0}^t\dd t_1\int_{t_0}^{t_1}\dd t_2\,
V_I(t_1)V_I(t_2)+\cdots.
\]
内层上限 \(t_1\) 正是时间排序：先发生的作用排在态的右侧。若有限维模型在整个时间窗内满足
\(\hbar^{-1}\int\|V_I(t)\|\dd t=\epsilon\ll1\)，将第 \(n\) 项的有序积分域嵌入完整的 \(n\) 维立方体，可得其算符范数不超过 \(\epsilon^n/n!\)。因此从二阶起的总振幅余项不超过 \(e^\epsilon-1-\epsilon\)。连续谱或无界相互作用时，不能机械使用这个范数界；那时需在所选波包和能窗上检验展开。

矩形时间窗 \(w(t)=1\)（\(-T/2<t<T/2\)）可直接积分：
\begin{align*}
\int_{-T/2}^{T/2}e^{\ii\Delta t}\dd t
&=\left.\frac{e^{\ii\Delta t}}{\ii\Delta}\right|_{-T/2}^{T/2}
=\frac{2\sin(\Delta T/2)}{\Delta}
=T\sinc(\Delta T/2).
\end{align*}
因此
\begin{equation}
P_{i\to f}^{(1)}(T)=\frac{|V_{fi}|^2T^2}{\hbar^2}
\sinc^2\!\left(\frac{\Delta_{fi}T}{2}\right),
\qquad \sinc x:=\frac{\sin x}{x},\quad\sinc0=1.
\label{eq:feqo-sinc-amplitude}
\end{equation}
首个零点在 \(|\Delta_{fi}|=2\pi/T\)。可分辨的能量尺度因此约为 \(2\pi\hbar/T\)；它由作用时间决定，不是初态能散，也不是谱仪的分辨率。时间窗若从 \(0\) 开始，振幅还乘 \(e^{\ii\Delta T/2}\)，概率不变，但相干叠加时不能丢掉这段相位。

对一般可积包络，定义 \(\widetilde w(\Delta)=\int w(t)e^{\ii\Delta t}\dd t\)。式~\eqref{eq:feqo-finite-time-amplitude} 表明谱线是 \(|\widetilde w|^2\)，并非所有脉冲都具有 sinc 旁瓣。例如 \(w(t)=e^{-t^2/(2\tau^2)}\) 给 \(\widetilde w=\sqrt{2\pi}\tau e^{-\Delta^2\tau^2/2}\)。这一例同时核对：时间宽度 \(\tau\) 越短，频率宽度越大。

这里容易混淆三种宽度。对矩形窗，\(\Delta=0\) 到首个零点的距离是 \(2\pi/T\)，但它不是概率峰的半高宽。对上述高斯窗，概率正比于 \(e^{-\Delta^2\tau^2}\)，半高点满足 \(|\Delta|=\sqrt{\ln2}/\tau\)，所以全半高宽为 \(2\sqrt{\ln2}/\tau\)。最后若入射电子本身有能量分布或谱仪有响应核，实验曲线还要对相应概率分布卷积。只从测得的峰宽倒推作用时间，必须先独立标定后两项。

\begin{proposition}[有限时间核的长时极限]
若 \(f(\Delta)\) 连续、有界且在零点附近变化缓慢，则
\begin{equation}
\lim_{T\to\infty}\int_{-\infty}^{\infty}\frac{\dd\Delta}{2\pi}
T\sinc^2(\Delta T/2)\,f(\Delta)=f(0).
\label{eq:feqo-finite-delta}
\end{equation}
\end{proposition}
\begin{proof}
令 \(x=\Delta T/2\)，积分变成
\(\int_{-\infty}^{\infty}(\dd x/\pi)\sinc^2x\,f(2x/T)\)。利用 \(\int_{-\infty}^{\infty}\sinc^2x\dd x=\pi\)，再由连续性与可积支配取极限即得。这里收敛的是与测试函数积分后的结果，\(T\sinc^2(\Delta T/2)\) 在每个非零 \(\Delta\) 处没有普通函数的极限。
\end{proof}

若末态在能量 \(E\) 上连续，\(\rho(E)\dd E\) 表示这段能量区间的末态数目。把 \(\Delta=(E-E_i)/\hbar-\omega\) 代入，注意 \(\dd E=\hbar\dd\Delta\)，便有
\begin{equation}
\frac{P(T)}{T}\longrightarrow
\frac{2\pi}{\hbar}|V(E_i+\hbar\omega)|^2\rho(E_i+\hbar\omega).
\label{eq:feqo-golden-rule-limit}
\end{equation}
这就是黄金律。它要求 \(|V(E)|^2\rho(E)\) 在宽度 \(\hbar/T\) 内几乎不变，同时一阶总概率仍远小于一；若共振跃迁已经反复发生，须回到式~\eqref{eq:feqo-interaction-evolution} 求解。

这两个条件并不矛盾，而是规定一个可能存在的中间时间窗。设末态连续谱的结构变化尺度为 \(E_c\)，黄金律要 \(T\gg\hbar/E_c\)；设式~\eqref{eq:feqo-golden-rule-limit} 的速率为 \(\Gamma\)，一阶理论又要 \(T\ll1/\Gamma\)。只有 \(\hbar/E_c\ll T\ll1/\Gamma\) 时，常速率与低总概率能够同时成立。若连续谱有很窄的共振或带边，\(E_c\) 变小，常速率窗口可能根本不存在。

一个离散二态反例说明为什么不可把黄金律用到任意长时间。若旋转绘景中的矩阵为
\[
H_{\rm rot}=\frac{\hbar}{2}
\begin{pmatrix}-\Delta&\Omega\\\Omega^*&\Delta\end{pmatrix},
\]
直接平方得 \(H_{\rm rot}^2=\hbar^2(\Delta^2+|\Omega|^2)I/4\)。把指数级数的偶次与奇次项分开，从第一态出发的第二态概率为
\[
P_{1\to2}(T)=\frac{|\Omega|^2}{\Delta^2+|\Omega|^2}
\sin^2\!\left(\frac{T}{2}\sqrt{\Delta^2+|\Omega|^2}\right).
\]
在 \(|\Omega|T\ll1\) 时它退化为 \(|\Omega|^2T^2\sinc^2(\Delta T/2)/4\)；再延长时间，概率却振荡而非无限线性增长。黄金律的长期线性增长来自\emph{连续末态的求和及合适的极限顺序}，不是单个离散跃迁的永久性质。

对动量明确的电子，交换 \((\hbar\omega,\hbar\mathbf q)\) 后还要检查电子色散。若 \(E(\mathbf p)=\sqrt{m^2c^4+c^2|\mathbf p|^2}\)，精确失谐为
\begin{equation}
\Delta(\mathbf p,\mathbf q)=\frac{E(\mathbf p+\hbar\mathbf q)-E(\mathbf p)}{\hbar}-\omega .
\label{eq:feqo-recoil-detuning}
\end{equation}
相互作用时间内 \(|\Delta|T\lesssim1\) 的通道能够获得显著振幅。小的 \(\hbar\omega/E\) 本身不足以保证低反冲：决定能否把多个边带视作等间隔的是这些边的失谐差乘以 \(T\)。

\begin{exercise}
求长度相同的矩形窗与高斯窗在失谐 \(\Delta=2/\tau\) 时的振幅，并说明高斯脉冲为什么没有矩形窗的严格零点。
\end{exercise}


短时谱线使我们知道电子能把多少能量交给材料，却还不知道这些能量流向何处。材料的推迟 Green 张量把同一电流激发的所有通道放在一个表达式里，因而可以在相同起点上比较电子损失与远场发光。

\section{一条电子电流怎样同时激发损失与发光}
\label{chap:feqo-eels-cl-kernel}

前章的 Green 张量回答给定电流会产生什么场。这里把电流指定为电子轨迹，并追踪电子交出的能量；对应实验分别称电子能量损失谱学（EELS）和阴极发光（CL）。先采用经典直线轨迹 \(\mathbf r_e(t)=\mathbf R_0+\mathbf vt\) 作为窄波包的中心；电子电荷始终记为 \(q=-e\)。相应的电流密度为
\begin{equation}
\mathbf j_e(\mathbf r,t)=q\mathbf v\,\delta^{(3)}(\mathbf r-\mathbf r_e(t)).
\label{eq:feqo-electron-current}
\end{equation}
电流的单位为 \(\si{\ampere\per\square\metre}\)。按全书 \(e^{-\ii\omega t}\) 约定，正频 Fourier 分量为 \(\mathbf j_e(\mathbf r,\omega)=\int_{-\infty}^{\infty}\mathbf j_e(\mathbf r,t)e^{\ii\omega t}\dd t\)。对 \(\mathbf v=v\hat{\mathbf z}\) 且 \(\mathbf R_0\) 固定，利用 \(\delta(z-vt)\) 的 Jacobian，得到
\[
\mathbf j_e(\mathbf r,\omega)=q\hat{\mathbf z}\,
\delta^{(2)}(\mathbf R-\mathbf R_0)e^{\ii\omega z/v}.
\]
这条沿 \(z\) 延伸的源说明电子读取的是一整段空间响应，而非轨迹上一点的局域态密度。

定义由样品引起的散射 Green 张量 \(\mathbf G_{\rm sc}=\mathbf G-\mathbf G_0\)，并令 \(\mathbf E_{\rm ind}=\ii\mu_0\omega\mathbf G_{\rm sc}\mathbf j_e\)，其中乘积含两点空间积分。电子对诱导场所做的功为 \(W_e=\int\dd t\,\mathbf j_e\cdot\mathbf E_{\rm ind}\)。损失定义为 \(W_{\rm loss}=-W_e\)。把实场写成正负频两项，频率 \(\omega>0\) 的谱密度为
\begin{equation}
\begin{aligned}
\frac{\dd W_{\rm loss}}{\dd\omega}
&=-\frac1\pi\operatorname{Re}\int\dd^3r\,
\mathbf j_e^*(\mathbf r,\omega)\cdot\mathbf E_{\rm ind}(\mathbf r,\omega)\\
&=\frac{\mu_0\omega}{\pi}
\langle\mathbf j_e,\operatorname{Im}\mathbf G_{\rm sc}\,\mathbf j_e\rangle .
\end{aligned}
\label{eq:feqo-eels-work}
\end{equation}
第一步用了实场的正负频配对及 Parseval 恒等式，第二步用了 \(\operatorname{Re}(\ii z)=-\operatorname{Im}z\)。二次型的两端各含一个电子电流，并在 \(\mathbf r,\mathbf r'\) 上积分。此处的虚部是算符虚部，定义为
\[
\operatorname{Im}\mathbf G
=\frac{\mathbf G-\mathbf G^\dagger}{2\ii}.
\]
对样品额外损失，减去真空项很重要：不能把自由空间核的奇异自场误算作材料信号。

将上面已经求出的直线电流代入二次型，二维横向 \(\delta\) 函数立即完成两个横向积分，剩下
\begin{equation}
\begin{aligned}
\frac{\dd W_{\rm loss}}{\dd\omega}
={}&\frac{\mu_0\omega q^2}{\pi}
\int\dd z\dd z'\,e^{\ii\omega(z'-z)/v}\\
&\times\operatorname{Im}G_{{\rm sc},zz}
\bigl((\mathbf R_0,z),(\mathbf R_0,z');\omega\bigr).
\end{aligned}
\label{eq:feqo-trajectory-green-kernel}
\end{equation}
这里的 \(\operatorname{Im}G_{zz}\) 是两点核的算符虚部：若环境不满足互易性，不能只把核的每个复数值逐点取虚部。\(z,z'\) 两处的场响应会乘电子沿途的相位差再干涉，因而 EELS 不只是某一点的局域光学态密度。这个双重积分也是平面界面算例的直接起点。

弱激发时，每个频率为 \(\omega\) 的独立量子携带 \(\hbar\omega\)。单次损失概率的频率密度可由能量谱定义为
\[
\Gamma_{\rm loss}(\omega)=\frac1{\hbar\omega}
\frac{\dd W_{\rm loss}}{\dd\omega}.
\]
这个等同需要单量子、稀事件极限；强激发时平均功仍由式~\eqref{eq:feqo-eels-work} 描述，却不能把平均量直接当成“恰好一次损失”的概率。

\begin{proposition}[被动环境中的正性]
若环境的完整响应被动，\(\langle\mathbf j,\operatorname{Im}\mathbf G\mathbf j\rangle\geq0\) 对任意可容许电流成立；因此总耗散功率非负。
\end{proposition}
\begin{proof}
驱动电流为场做的周期平均功率是 \((1/2)\operatorname{Re}\langle\mathbf j,\mathbf E\rangle\)。代入 \(\mathbf E=\ii\mu_0\omega\mathbf G\mathbf j\)，并按吸收功率的相反符号取正，得到 \((\mu_0\omega/2)\langle\mathbf j,\operatorname{Im}\mathbf G\mathbf j\rangle\)。被动性正是吸收功率对任意源非负。散射部分单独不一定正；它表示相对真空的增减，不能对它无条件使用本命题。
\end{proof}

发光测量问同一诱导场有多少流过远处的收集面。令 \(\mathbf E_{\rm far}(r\hat{\mathbf n},\omega)=e^{\ii kr}\mathbf F(\hat{\mathbf n},\omega)/r+O(r^{-2})\)。真空远场中 \(\mathbf H_{\rm far}=\hat{\mathbf n}\times\mathbf E_{\rm far}/Z_0\)，其中 \(Z_0=\sqrt{\mu_0/\varepsilon_0}\) 为真空波阻抗。因 \(\mathbf F\) 垂直于 \(\hat{\mathbf n}\)，频域 Poynting 通量乘面积元 \(r^2\dd\Omega\) 后，角向辐射能量是 \(\dd W_{\rm rad}/(\dd\omega\dd\Omega)=|\mathbf F|^2/(\pi Z_0)\)；这里的 \(1/\pi\) 与正频 Fourier 约定相配。探测器只取方向集合 \(\Omega_{\rm det}\) 及偏振 \(\lambda\)，其理想光子产额为
\begin{equation}
\Gamma_{\rm CL}(\omega)=\frac1{\hbar\omega}
\sum_\lambda\int_{\Omega_{\rm det}}\dd\Omega\,
\eta(\omega,\hat{\mathbf n},\lambda)
\frac{\dd W_{{\rm rad},\lambda}}{\dd\omega\dd\Omega}.
\label{eq:feqo-cl-collection}
\end{equation}
\(\eta\in[0,1]\) 是收集与探测效率。同一个 \(\mathbf G\mathbf j_e\) 先给总诱导场，再给远场 \(\mathbf F\)。材料吸收、导波和未被收集的辐射都可能增加电子损失，却不会增加此探测器的计数。

这个区别还能用一个单共振响应直接核对。令复模式振幅 \(a(t)\) 的模平方等于储能（单位 \(\si{\joule}\)），并用线性方程 \(\dot a=(-\ii\omega_0-\kappa/2)a+s(t)\) 描述电子电流的驱动；这里 \(s\) 的单位是 \(\si{\joule}^{1/2}\si{\per\second}\)。辐射与材料吸收的能量衰减率分别为 \(\kappa_r,\kappa_a\)，总宽度 \(\kappa=\kappa_r+\kappa_a\)。Fourier 变换给 \(|a(\omega)|^2=|s(\omega)|^2/[(\omega-\omega_0)^2+(\kappa/2)^2]\)。令 \(A(\omega)=|s(\omega)|^2/(2\pi)\)，它的单位是 \(\si{\joule}\)；假设脉冲前后模式储能都回到零，复振幅的全频 Parseval 恒等式便把电子总损失、辐射能量和收集能量的频谱写为
\[
\frac{A(\omega)\kappa}{(\omega-\omega_0)^2+(\kappa/2)^2},\qquad
\frac{A(\omega)\kappa_r}{(\omega-\omega_0)^2+(\kappa/2)^2},\qquad
\eta_{\rm coll}(\omega)
\frac{A(\omega)\kappa_r}{(\omega-\omega_0)^2+(\kappa/2)^2}.
\]
三个表达式的单位都是 \(\si{\joule\second}\)，正好是 \(\dd W/\dd\omega\) 的单位。这不是对任意材料的定理，而是一个清楚标明单模、线性响应与无其他通道假设的算例。它告诉我们：增加材料吸收率会改变总峰宽，也会降低辐射占比；扩大收集角只能改变最后一个量。各谱线看似有相同峰位，并不意味着 EELS 与 CL 的峰高可以直接相除当作仪器效率。

真空中无限长的匀速轨迹提供一个必要检查。轨迹的 Fourier 支持为 \(k_z=\omega/v\)，自由光子的壳条件为 \(|\mathbf k|=\omega/c\)。因为 \(v<c\) 使 \(|k_z|>\omega/c\)，两者不相交，所以没有真实远场辐射。界面或周期结构能提供额外动量，才可能改变这个结论。

\begin{exercise}
用 \(\delta(z-vt)=\delta(t-z/v)/v\) 重新推导电子电流的频域表达，并核对其量纲。
\end{exercise}


一般环境含连续多个通道。若实验只在一个窄的频率与动量窗口工作，才可把它投影到少数通道。这个步骤保留电子的精确色散，不能同时预设等间隔边带。

\section{保留色散的离散反冲格}
\label{chap:feqo-recoil-lattice}

若近场的一个 Fourier 分量每次给电子动量 \(\hbar\mathbf q\)，从初始动量 \(\mathbf p_0\) 可达的态便标为整数 \(\ell\)：\(\mathbf p_\ell=\mathbf p_0+\ell\hbar\mathbf q\)。这个标号只记录交换次数；它不宣称电子的真实能量间隔相等。各态的精确能量为
\begin{equation}
\mathcal E_\ell=\sqrt{m^2c^4+c^2|\mathbf p_0+\ell\hbar\mathbf q|^2},\qquad
\delta_\ell=\mathcal E_\ell-\mathcal E_0-\ell\hbar\omega .
\label{eq:feqo-channel-lattice}
\end{equation}
\(\delta_\ell\) 是以 \(\ell\hbar\omega\) 为参考的能量失谐，\(\delta_0=0\)。相邻跃迁的频率失谐是 \((\delta_{\ell+1}-\delta_\ell)/\hbar\)，恰好回到式~\eqref{eq:feqo-recoil-detuning}。

把电子波函数写成 \(\sum_\ell c_\ell(t)e^{-\ii(\mathcal E_0+\ell\hbar\omega)t/\hbar}|\mathbf p_\ell\rangle\)，再将周期势的正负频项代入薛定谔方程，按 \(|\mathbf p_\ell\rangle\) 投影，得到
\begin{equation}
\ii\dot c_\ell=\frac{\delta_\ell}{\hbar}c_\ell
+g_{\ell-1}c_{\ell-1}+g_\ell^*c_{\ell+1}.
\label{eq:feqo-exact-recoil-chain}
\end{equation}
\(g_\ell\) 是从 \(\ell\) 到 \(\ell+1\) 的角频率量纲耦合；厄米 共轭使反向边的系数为共轭，因此截断边界没有概率流出时 \(\sum_\ell|c_\ell|^2\) 守恒。式~\eqref{eq:feqo-exact-recoil-chain} 在给定单频单动量模式投影内是精确的离散方程；若场含连续频率或动量，须回到原来的积分方程。

一维纵向交换使物理尺度更清楚。取 \(\mathbf q=k_z\hat{\mathbf z}\)，其中 \(k_z\) 是带符号的纵向波数，不能与电子电荷 \(q=-e\) 混淆。设 \(p_0>0\)、\(v_0=\partial_pE(p_0)\)、\(\gamma_0=E_0/(mc^2)\)。Taylor 定理给
\begin{align*}
\mathcal E_\ell
&=E_0+\ell\hbar k_zv_0
+\frac{(\ell\hbar k_z)^2}{2\gamma_0^3m}+R_{3,\ell},\\
|R_{3,\ell}|&\leq\frac{|\ell\hbar k_z|^3}{6}
\max_{|p-p_0|\leq|\ell\hbar k_z|}|E'''(p)|.
\end{align*}
定义反冲角频率 \(\omega_r:=\hbar k_z^2/(2\gamma_0^3m)\)。于是
\begin{equation}
\frac{\delta_\ell}{\hbar}
=\ell(k_zv_0-\omega)+\ell^2\omega_r+\frac{R_{3,\ell}}{\hbar}.
\label{eq:feqo-recoil-expansion}
\end{equation}
首项是相位速度失配；二次项使不同边不能同时严格共振。若相互作用实际占据 \(|\ell|\leq N\)，低反冲近似需同时满足 \(N|k_zv_0-\omega|T\ll1\)、\(N^2\omega_rT\ll1\) 和 \(T\max_{|\ell|\leq N}|R_{3,\ell}|/\hbar\ll1\)。此外，若用常数耦合，还需 \(T\max|g_\ell-g_0|\ll1\)。这些量控制的是相位误差，可随作用时间或边带数增加而失效。

沿同一格子可做一个贯穿下一章的共振选择算例。暂时保留二阶色散，并把驱动频率设为 \(\omega_d=k_zv_0+\omega_r\)，即不只匹配电子中心速度，还补上从 \(0\) 到 \(1\) 的反冲能。式~\eqref{eq:feqo-recoil-expansion} 此时变成 \(\delta_\ell/\hbar=\omega_r\ell(\ell-1)\)。\(\ell=0,1\) 两态同失谐，而最近的 \(-1,2\) 两态都相隔 \(2\omega_r\)。若 \(|g|/(2\omega_r)\ll1\)、\(|R_{3,\ell}|T/\hbar\ll1\) 且耦合时间不长到累积高阶能级位移，最低阶受限方程给从 \(|0\rangle\) 到 \(|1\rangle\) 的概率 \(\sin^2(|g|T)\)。被排除的两条边的振幅尺度是 \(O(|g|/\omega_r)\)，人口尺度为它的平方；下一章将从精确投影消元求出它们造成的能级修正。若仍沿用 \(\omega=k_zv_0\)，则 \(\delta_{\pm1}=\hbar\omega_r\)，连第一条边也不是严格共振。这说明速度同步与相邻能级共振在可分辨反冲时是两回事。

当这些条件成立，式~\eqref{eq:feqo-exact-recoil-chain} 中的对角失谐可在共同相位下略去，耦合近似与 \(\ell\) 无关；下一章的 Bessel 能量梳便从这个平移不变的格子出现。反过来，当只有某一条边近共振时，格子的其余通道必须定量消去，而非直接删掉。

\begin{exercise}
从 \(E(p)=\sqrt{m^2c^4+c^2p^2}\) 求出 \(E'(p_0)\) 与 \(E''(p_0)\)，再计算相邻两条边的失谐差至 \(k_z^2\) 阶。
\end{exercise}

